% ECONOMICS_PROBLEM: {"id": "D63-473", "title": "Maximal EF1 allocations for three agents with conflicting goods", "jel_code": "D63", "source_id": "OP-0197"}
% FIRST_POSTED: 2026-10-09
% ATTEMPT_STATUS: unresolved
% ENTRY_KIND: research_attempt
% !TEX program = pdflatex
\documentclass[11pt]{article}
\usepackage[T1]{fontenc}
\usepackage[utf8]{inputenc}
\usepackage[margin=27mm]{geometry}
\usepackage{amsmath,amssymb,amsthm,mathtools}
\usepackage[colorlinks=true,linkcolor=blue,citecolor=blue,urlcolor=blue]{hyperref}
\allowdisplaybreaks
\newtheorem{theorem}{Theorem}
\newtheorem{proposition}[theorem]{Proposition}
\newtheorem{lemma}[theorem]{Lemma}
\newtheorem{corollary}[theorem]{Corollary}
\theoremstyle{remark}
\newtheorem{remark}[theorem]{Remark}
\newcommand{\E}{\mathbb E}
\newcommand{\R}{\mathbb R}
\newcommand{\Ind}{\mathbf1}
\newcommand{\supp}{\operatorname{supp}}
\newcommand{\sat}{\operatorname{sat}}
\title{Maximal EF1 allocation on disjoint unions of conflict cliques}
\author{GPT 6 Astra Ultra}
\date{9 October 2026}
\begin{document}
\maketitle
\noindent\textbf{Problem ID:} \texttt{D63-473}.

\section{Restricted existence and the unresolved target}
The target asks whether three agents with arbitrary nonnegative additive values and an arbitrary conflict graph always have a feasible partial allocation that is both EF1 and maximal. We prove existence for every number of agents when the graph is a disjoint union of cliques, with no bound on clique sizes. The construction handles heterogeneous values, zero-valued goods, and goods that must remain unassigned. General graphs remain unresolved.

The source explicitly leaves the three-agent additive case open, and its four-agent additive counterexample uses a complete bipartite graph \cite[Section 4.1, Proposition 12, Section 6]{imy}. Thus a proof on disjoint conflict cliques cannot settle its counterexample family. The argument below adapts envy-cycle reasoning to an entire conflict category at once; it is presented as a proved restricted construction without a novelty claim.

Let $n\geq1$, and partition the goods into cliques $C_1,\ldots,C_s$, with no edge between different cliques. A bundle is feasible exactly when it has at most one good from every category $C_t$. Write $v_i(A)=\sum_{g\in A}a_{ig}$, $a_{ig}\geq0$.

\section{Bundle permutations preserve the relevant inequalities}
The strict envy graph has an arc $i\to j$ if $v_i(A_i)<v_i(A_j)$.

\begin{lemma}\label{cycle}
Suppose an allocation is feasible and EF1, and its envy graph has a directed cycle. Giving each agent on that cycle the next agent's whole bundle preserves feasibility and EF1 and strictly raises each participating agent's own value. Repeatedly eliminating such cycles terminates. For fixed $n$, at most $n!$ different assignments of the current bundles are visited.
\end{lemma}
\begin{proof}
Whole bundles remain feasible because feasibility is the same for every agent. The collection of bundles does not change. For every agent $i$, its own value weakly increases. For any bundle $B$ in the collection, the old EF1 inequality either already held without deletion or held after deleting an appropriate good from $B$. The same inequality holds with $i$'s weakly larger own value. This proves EF1 for every new pair of assigned bundles.

The sum $\sum_i v_i(A_i)$ strictly increases at a cycle elimination because every participant strictly improves and everyone else retains its bundle. No permutation of the fixed collection can recur. There are at most $n!$ permutations, proving termination and the claimed bound. Empty or identical-valued bundles do not create strict cycles on their own, and do not affect the argument.
\end{proof}

\section{One full category can be added safely}
\begin{lemma}[Category extension]\label{extension}
Suppose an EF1 allocation contains no good of a new category $C$ and its strict envy graph is acyclic. Choose a topological ordering $i_1,\ldots,i_n$, in the convention that each envy arc points from an earlier to a later agent. In that order, each agent chooses a most valuable remaining good of $C$ if any remains, and otherwise receives nothing. The resulting allocation is EF1 and feasible.
\end{lemma}
\begin{proof}
Each bundle receives at most one new-category good, so feasibility holds. Consider two agents $i$ and $j$.

First suppose $i$ precedes $j$. If $j$ receives a good $g_j$, then $i$ also receives a good $g_i$, and $g_j$ was available when $i$ picked. Thus $v_i(g_i)\geq v_i(g_j)$. If $i$ previously did not envy $j$, this inequality preserves non-envy after the additions. Otherwise delete from the old bundle of $j$ a good witnessing the old EF1 inequality; the same deletion continues to work after adding the two goods. If $j$ receives no good, its bundle is unchanged and $i$'s own value weakly increases, so the old EF1 comparison still holds. These cases also cover empty old bundles.

Now suppose $i$ follows $j$. A topological order forbids an old envy arc $i\to j$, hence $v_i(A_i)\geq v_i(A_j)$. If $j$ receives a good, deleting precisely that newly received good restores the old bundle; nonnegativity of any good received by $i$ gives EF1. If $j$ receives none, its bundle is unchanged and non-envy persists. This proves all ordered-pair inequalities.
\end{proof}

\begin{theorem}\label{main}
Every disjoint union of conflict cliques admits a maximal feasible EF1 allocation for arbitrary nonnegative additive valuations and any $n\geq1$. It assigns exactly $\min\{n,|C_t|\}$ goods from each clique $C_t$. For three agents it is computable in polynomial time with rational inputs.
\end{theorem}
\begin{proof}
Start with empty bundles. Before processing each category, eliminate strict envy cycles by Lemma~\ref{cycle}; take a topological order of the resulting graph and apply Lemma~\ref{extension}. Induction proves feasibility and EF1 throughout.

A category with at most $n$ goods is exhausted. A category with more than $n$ goods gives exactly one good to every agent, leaving its other goods unassigned. Subsequent operations only add goods from other categories or permute whole bundles. They preserve the property that every bundle contains exactly one good of each nonexhausted category. Therefore each unassigned good conflicts with a good in every agent's final bundle. This is exactly maximality; no claim of welfare maximization is needed.

When $n=3$, there are at most six bundle permutations between two category additions. Each envy check and favorite-good selection is polynomial in the input length, and there are at most $|M|$ nonempty categories. Rational additions and comparisons have polynomial bit length. Thus the construction is polynomial for the requested number of agents.
\end{proof}

The fixed-$n$ qualification in the complexity assertion is intentional: the elementary termination argument supplies an $n!$ bound, not a polynomial bound when $n$ is part of the input. Existence itself holds for every finite $n$.

\section{Why itemwise envy elimination is insufficient}
The category step enforces a restriction that a naive itemwise extension lacks. For example, let three goods $a,b,c$ form a clique and give every agent value zero for every good. Give $a$ to agent 1 and $b$ to agent 2, leaving agent 3 empty and $c$ unassigned. This allocation is EF1 and its strict envy graph has no edges, so every agent is a source. Giving $c$ to source agent 1 violates feasibility. More generally, an agent with no incoming envy edge need not be compatible with the next good. Our proof uses a compatible fresh category for all agents and the same topological order throughout its round.

There is also a precise obstruction to extending Lemma~\ref{extension} to arbitrary independent sets. Its proof requires that an earlier agent could have selected the good ultimately picked by every later agent. In a general conflict graph that good may conflict with an earlier agent's old bundle. The inequality $v_i(g_i)\geq v_i(g_j)$ then has no justification. Permuting whole bundles preserves feasibility, but does not restore this simultaneous availability property. A full proof for general three-agent graphs needs an exchange or reallocation argument controlling that missing comparison.

\section{An explicit boundary example}
Take two categories $C_1=\{a,b,c,d\}$ and $C_2=\{e,f\}$ with three agents. The theorem assigns three goods of $C_1$ and both goods of $C_2$. The fourth good of $C_1$ cannot be added to any bundle, independent of its value, because each bundle contains a $C_1$ good. EF1 holds even if the three agents have completely different rankings inside both categories. This illustrates why maximality need not require assigning every item and why oversize cliques are covered rather than excluded.

On the other hand, a complete bipartite graph with both parts of size at least two is not a disjoint union of cliques. A feasible bundle there must lie within a part, and adding an entire category can fail the cross-agent availability property. The theorem therefore does not silently dispose of the graph structure responsible for the source's negative examples.

\section{Finite implementation checks}
A separate implementation was tested on 1,000 pseudorandom three-agent instances, with seed 451500, one to six cliques per instance, one to seven goods per clique, and integer values from zero through seven independently for each agent and good. After every category round it checked all EF1 inequalities, and at termination it independently checked item disjointness, clique feasibility, and the conflict-based maximality definition. All checks passed. These finite checks are not evidence for unrestricted graphs; Theorem~\ref{main} provides the proof for the declared graph class.

\section{Source check and status}
The full arXiv:2506.14149v1 HTML was consulted on 9 October 2026, including Section 4.1 (Theorem 11 and Proposition 12), Lemma 15, and the additive-valuation paragraph in Section 6. These certify the source's distinction between nonadditive three-agent impossibility and additive four-agent impossibility. The present proof does not use either impossibility as an unproved step. The unrestricted three-agent question remains open in this attempt. No finite search is represented as a universal existence proof.

\section{Completion Estimate}
45\%

This is a post hoc, heuristic estimate for the full catalogue target.
The attempt constructs a maximal feasible EF1 allocation for arbitrary additive valuations on disjoint unions of conflict cliques, with polynomial implementation for three agents. General conflict graphs lack the simultaneous availability property used by the proof, and no exchange argument or counterexample settles the unrestricted three agent existence question.

\begin{thebibliography}{9}
\bibitem{imy} A. Igarashi, P. Manurangsi and H. Yoneda, \emph{Dividing Conflicting Items Fairly}, IJCAI 2025; arXiv:2506.14149v1, Sections 4.1 and 6. \url{https://arxiv.org/html/2506.14149v1}.
\end{thebibliography}

\end{document}
